\documentclass[12pt]{article}
\usepackage[english]{babel}
\usepackage[utf8]{inputenc}
\usepackage{graphicx}
\usepackage{amsmath}
\usepackage{amsfonts}
\usepackage{amssymb}
\usepackage{float}
\usepackage{hyperref}
\usepackage{geometry}
\usepackage{titling}
\geometry{a4paper, total={170mm,257mm}, left=20mm, top=20mm}

\title{The Power of Preprocessing in Image Classification: A Case Study with Breast Cancer Images}
\author{Pedro Figueira Bôa-Viagem}
\date{\small July 22, 2025}

\renewcommand{\maketitle}{
  \begin{flushleft}
    {\LARGE\bfseries \thetitle}
    \vspace{1em}
    
    \begin{minipage}[t]{0.5\textwidth}
      \raggedright
      \theauthor
    \end{minipage}%
    \begin{minipage}[t]{0.5\textwidth}
      \raggedleft
      \thedate
    \end{minipage}
  \end{flushleft}
  \vspace{2em}
}

\begin{document}

\maketitle

\section*{Abstract}
Leveraging machine learning techniques for image classification, particularly in the context of breast cancer detection, requires meticulous data preparation and a robust dataset to ensure accurate model performance. This document outlines the data augmentation and preprocessing techniques applied to a breast cancer image dataset, enhancing the analysed models ability to generalize and perform well in real-world scenarios. The methods described here are essential for improving the robustness of machine learning models against common variations in image data and ensuring that the models can effectively learn from the available data with help of simple image processing techniques.

TODO - expand (results, methods, models, preprocessing, etc.)

\section*{1. Introduction}

Breast cancer is one of the most prevalent cancers worldwide, and early detection is crucial for improving patient outcomes. Medical imaging, particularly mammography, plays a central role in screening and diagnosis. However, interpreting these images is a challenging task, often requiring expert radiologists and being subject to inter-observer variability.

In recent years, machine learning and deep learning techniques have shown great promise in automating image classification tasks, including the detection and grading of breast cancer from mammograms. Despite their potential, the performance of these models is highly dependent on the quality and diversity of the training data. Variations in image acquisition, noise, and limited dataset size can hinder model generalization and robustness.

To address these challenges, effective data preprocessing strategies are essential. Preprocessing techniques such as denoising, binarization, and morphological operations can enhance image quality and highlight relevant features, while data augmentation increases dataset diversity by simulating real-world variations. Together, these steps help build more accurate and reliable models for breast cancer image classification.

This work presents a systematic study of preprocessing techniques applied to a breast cancer image dataset. We evaluate the impact of different preprocessing pipelines on the performance of several deep learning models, aiming to identify the most effective strategies for improving classification accuracy in this critical medical application.

This document is structured as follows: Section 2 describes the data augmentation techniques applied to the dataset, as well as the data augmentation for the data. Then Section 3 details the image preprocessing steps, and Section 4 presents the results of the experiments conducted with various machine learning models. Finally, Section 5 concludes with a discussion of the findings and their implications for future research in this area.

\section*{2. Dataset Description and Augmentation}

\subsection*{2.1 Dataset Description}

The dataset used in this analysis is the \textit{CBIS-DDSM: Breast Cancer Image Dataset}, which is a collection of mammography images annotated with BI-RADS (Breast Imaging Reporting and Data System) labels. This dataset is widely used in the field of medical imaging and provides a rich source of data for training and evaluating machine learning models. Although the dataset is relatively small, with only 410 images, it is well-suited for this case study due to its high-quality annotations and the presence of various breast cancer cases. It will also be augmented to increase the effective dataset size and diversity.

\subsection*{2.2. Data Augmentation and Selection}

In this section, we describe the techniques applied to prepare the images in the database before their use in the machine learning model. The steps follow modern practices of data augmentation and image normalization, in order to improve the model's generalization and increase robustness against common variations in real-world environments.

The original mammography images, provided in DICOM format, were first converted to PNG and normalized to the [0, 1] range for better model understanding. To address the limited size and variability of the dataset, we applied a comprehensive data augmentation pipeline using the \texttt{albumentations} library. The following random transformations were used to simulate real-world variations and increase the diversity of the dataset:

\begin{itemize}
    \item \textbf{Horizontal and vertical flips}: simulate different orientations of the breast tissue;
    \item \textbf{Random brightness and contrast adjustment}: account for differences in image acquisition conditions;
    \item \textbf{Rotation}: random rotations up to 360 degrees to simulate different patient positions;
    \item \textbf{Affine transformations}: scaling and rotation to introduce geometric variability.
\end{itemize}

Each original image was augmented once for each type, resulting in several synthetic variants per sample. The augmentation process was implemented as follows:

\begin{verbatim}
import albumentations as A

# augmentation pipeline (randomly flip, brightness, rotation)
augment = A.Compose([
    A.HorizontalFlip(p=1.0),
    A.VerticalFlip(p=1.0),
    A.RandomBrightnessContrast(brightness_limit=0.1, contrast_limit=0.1, p=1.0),
    A.Rotate(limit=360, p=1.0, border_mode=cv2.BORDER_REPLICATE),
    A.Affine(scale=(0.6, 1.4), rotate=(-90, 90), p=1.0)
])
\end{verbatim}

For each image, the five new augmented versions were generated and saved, in addition to the original. All images were then resized to $224 \times 224$ pixels to ensure compatibility with standard deep learning architectures.

After augmentation, it was observed that the distribution of BI-RADS labels in the dataset was highly imbalanced, with some classes significantly underrepresented. To ensure fair model training and evaluation, we applied a balancing procedure by limiting the number of samples per class to a fixed maximum (e.g., 35 images per class). This trimming step reduced the total number of images from 1640 to 274, with each class represented equally. The balanced dataset was then split into training and test sets using a 75/25 ratio, maintaining the class distribution (stratified split). Only images with valid BI-RADS labels (1, 2, 3, 4a, 4b, 4c, 5, 6) were retained for further analysis (no images were discarded).

This augmentation and selection strategy significantly increased the effectiveness of the dataset diversity, with a final count of 274 of vastly different images, helping to prevent overfitting and improve the robustness of the trained models. The resulting dataset was used as input for the subsequent preprocessing and model training steps described in the following sections.

\section*{3. Image Preprocessing}

In this study, we selected a set of fundamental image preprocessing algorithms that are simple and widely used in medical image analysis and can be efficiently automated. The goal was to evaluate the impact of each transformation, both individually and in combination, on the performance of deep learning models for breast cancer classification. The chosen techniques, implemented using OpenCV and NumPy in Python, are as follows:

\begin{itemize}
    \item \textbf{Denoising (Gaussian Blur):} Reduces random noise in the images, making relevant structures more distinguishable and improving model robustness.
    \item \textbf{Binarization (Thresholding):} Converts grayscale images to binary by applying a fixed threshold, which can highlight regions of interest and simplify feature extraction.
    \item \textbf{Low-pass Filtering (Averaging):} Smooths the image by averaging pixel values in a local neighborhood, further reducing noise and minor variations.
    \item \textbf{Morphological Operations:} Includes erosion, dilation, opening (erosion followed by dilation), and closing (dilation followed by dilation). These operations are used to remove small artifacts, fill gaps, and enhance the shape of anatomical structures in the mammograms.
    \item \textbf{Contrast Limited Adaptive Histogram Equalization (CLAHE):} Enhances local contrast in the image, making it easier to distinguish features in regions with varying illumination.
\end{itemize}

Each preprocessing function was applied to the dataset as a separate experiment, with varying parameters. This systematic approach allowed us to identify which preprocessing strategies most effectively enhance model performance, while keeping the computational requirements low and the pipeline easily reproducible.

The implementation of these transformations is modular, enabling straightforward integration into the training workflow. This design ensures that the preprocessing steps can be adapted or extended for future studies or different datasets with minimal effort.

\begin{table}[H]
\centering
\begin{tabular}{|l|l|l|}
\hline
\textbf{Method} & \textbf{Parameter} & \textbf{Values Tested} \\
\hline
Gaussian Blur & Kernel Size & 3, 5, 7 \\
              & Sigma       & 0.5, 1.0, 2.0 \\
\hline
Thresholding  & Threshold Value & 80, 120, 160 \\
\hline
Averaging     & Kernel Size & 3, 5, 7 \\
\hline
Morphological & Operation & Erosion, Dilation, Opening, Closing \\
              & Kernel Size & 3, 5 \\
\hline
CLAHE         & Clip Limit & 2.0, 4.0 \\
              & Tile Grid Size & (8,8), (16,16) \\
\hline
\end{tabular}
\caption{Parameter variations tested for each preprocessing method.}
\label{tab:preprocessing_params}
\end{table}

\section*{4. Model Selection and Training}

To systematically evaluate the impact of each preprocessing strategy, we trained and validated several deep learning models on every preprocessed version of the dataset. The workflow was implemented in Python using Keras and TensorFlow, and consisted of the following steps:

\begin{itemize}
    \item \textbf{Preprocessing Application:} Each image in the training and test sets was processed using one of the defined preprocessing functions (denoising, binarization, low-pass filtering, or morphological operations), as well as all possible pairs of these operations in both sequential orders. This resulted in multiple distinct datasets, each corresponding to a unique preprocessing pipeline.
    \item \textbf{Model Architectures:} For each preprocessed dataset, three types of neural network models were trained:
          \begin{itemize}
              \item A custom convolutional neural network (CNN) with three convolutional layers, max pooling, and dense layers.
              \item A transfer learning model based on ResNet50, with the convolutional base frozen and custom dense layers added for classification.
              \item A transfer learning model based on DenseNet121, similarly adapted for the task.
          \end{itemize}
    \item \textbf{Training Procedure:} Each model was trained for up to 15 epochs using the Adam optimizer and categorical cross-entropy loss. Early stopping was not used to ensure comparability across runs. For transfer learning models, grayscale images were repeated across three channels to match the expected input shape.
    \item \textbf{Evaluation:} After training, the best validation accuracy achieved for each model and preprocessing combination was recorded, along with the corresponding epoch. This allowed for a direct comparison of the effectiveness of each preprocessing strategy and model architecture.
\end{itemize}

The entire process was automated to ensure reproducibility and fairness in comparison. The results, including the best validation accuracies for all combinations, are summarized in the next section. This systematic approach enabled us to identify which preprocessing techniques and model architectures yielded the highest classification performance on the breast cancer image dataset.

\section*{5. Results and Discussion}

TODO

\section*{6. Conclusion}

TODO

\end{document}